% Annex to capability-proof-arxiv.tex — included via \input{}.
\appendix
\section{Annex: one paper, end to end}
\label{sec:annex}

The warrant table asserts that the stages work. This annex shows one paper's
artifacts against its own source, so the assertion can be checked rather than
believed. The paper is \texttt{arXiv:2311.05789}, chosen because it is small
enough to display whole and complete across every layer: 2 proof passages, 2
CLean typings, 6 expository scopes. Nothing here is reformatted for
presentation --- the EDN is as the pipeline emitted it, and the source lines
are as \texttt{dp\_paper\_view} numbers them.

\subsection*{A.1 The source, with what each stage took from it}

\begin{center}
\small
\begin{tabular}{r|p{0.52\textwidth}|p{0.29\textwidth}}
\textbf{line} & \textbf{source (\LaTeX, truncated)} & \textbf{stage annotation}\\
\hline
\multicolumn{3}{l}{\emph{\quad Introduction --- expository region, S4 window L102--120}}\\
\hline
104 & \texttt{Monoidal categories were defined in \textbackslash cite\{be\} and further studied in \dots} & S1 citation marks\\
112 & \texttt{The concept of \{\textbackslash em fusion\} \dots\ originated in high energy physics} & \\
114 & \texttt{Currently the largest consumer of fusion categories is condensed matter physics (topological states of matter)} & \textbf{S4} $\to$ \texttt{:connection/}\linebreak\texttt{application-domain}\\
\hline
\multicolumn{3}{l}{\emph{\quad Definitions --- proof passage \texttt{p0}, S3 window L170--180}}\\
\hline
174 & \texttt{commutes for any \$X,Y\textbackslash in\textbackslash C\$.} & \\
176 & \texttt{Clearly composites of semi-groupal functors and natural transformations are semi-groupal.} & \textbf{S3} node \texttt{:semi-groupal-}\linebreak\texttt{composition} (\emph{anchored} \texttt{171--172}; \textbf{H21})\\
179 & \texttt{A semi-groupal category \$\textbackslash C\$ is \{\textbackslash em monoidal\} if it has an object \$I\textbackslash in\textbackslash C\$ \dots} & \textbf{S3} node \texttt{:monoidal-}\linebreak\texttt{category-def} (\emph{anchored} \texttt{173--180})\\
180 & \texttt{\textbackslash beq\textbackslash lb\{uc\}\textbackslash xygraph\{} & \textbf{S3} node \texttt{:unit-}\linebreak\texttt{coherence-diagram}\\
\hline
\multicolumn{3}{l}{\emph{\quad Module categories --- proof passage \texttt{p1}, S3 window L307--331}}\\
\hline
307 & \texttt{A \{\textbackslash em pivotal structure\} on a rigid monoidal \$\textbackslash C\$ is a monoidal isomorphism \$\textbackslash tau\$ \dots} & \textbf{S3} \texttt{:pivotal-def} \emph{(307--307 --- exact)}\\
308 & \texttt{The group \$Aut\_\textbackslash ot(Id\_\textbackslash C)\$ acts freely on the set of pivotal structures \dots} & \textbf{S3} \texttt{:aut-action} \emph{(308--308 --- exact)}\\
310 & \texttt{For an endomorphism \$f:X\textbackslash to X\$ \dots\ \$tr(f)=ev\_\{X\^{}*\}(\textbackslash tau\_X f\textbackslash ot 1)coev\_X\$} & \textbf{S3} \texttt{:trace-def} \emph{(310--310 --- exact)}\\
311 & \texttt{A pivotal category \$\textbackslash C\$ is \{\textbackslash em spherical\} if \$tr(f\^{}*) = tr(f)\$ \dots} & \textbf{S3} \texttt{:spherical-def} \emph{(311--311 --- exact)}\\
319 & \texttt{A category \$\textbackslash M\$ is a \{\textbackslash em (left) module\} category over \$\textbackslash C\$ \dots} & \textbf{S3} \texttt{:module-}\linebreak\texttt{category-def} \emph{(exact)}\\
325 & \texttt{A natural transformation \$c:F\textbackslash to G\$ between \$\textbackslash C\$-module functors \dots} & \textbf{S3} \texttt{:module-natural-}\linebreak\texttt{trans} \emph{(anchored 327 --- off by 2)}\\
\hline
\end{tabular}
\end{center}

Read the third column as the pipeline's own claim about the second. The
\emph{content} is faithful throughout: every node text is a recognisable gloss
of the line it names, and none is invented. The \emph{anchors} are exact for
five of six nodes in \texttt{p1} and drift four to five lines in \texttt{p0} ---
which is \hz{H21}, visible here in a way no aggregate count reached, and the
reason this annex exists.

\subsection*{A.2 What S3 produced (verbatim, passage \texttt{p0})}

\begin{center}\small
\begin{minipage}{0.93\textwidth}\ttfamily\raggedright\footnotesize
\{:paper/id "2311.05789", :passage/id "2311.05789:monoidal-def:170-180",\\
\ :source \{:lines [170 180], :kind :proof\},\\
\ :nodes [\{:id :semi-groupal-composition, :kind :claim,\\
\ \ \ \ \ \ \ \ \ :text "composites of semi-groupal functors and natural transformations are semi-groupal",\\
\ \ \ \ \ \ \ \ \ :source \{:lines [171 172]\}\}\\
\ \ \ \ \ \ \ \ \{:id :monoidal-category-def, :kind :claim, :text "a semi-groupal category C is monoidal if it\\
\ \ \ \ \ \ \ \ \ has a monoidal unit I and natural isomorphisms l\_X, r\_X satisfying the unit coherence condition"\dots\}\\
\ \ \ \ \ \ \ \ \{:id :unit-coherence-diagram, :kind :object, \dots\}],\\
\ :edges [\{:id :e-monoidal-def, :kind :infer, :relation :defines,\\
\ \ \ \ \ \ \ \ \ :premise :semi-groupal-composition, :conclusion :monoidal-category-def,\\
\ \ \ \ \ \ \ \ \ :warrant \{:kind :missing-warrant,\\
\ \ \ \ \ \ \ \ \ \ \ \ \ \ \ \ \ :wanted :semi-groupal-composition-is-closed-under-composition\}\}\\
\ \ \ \ \ \ \ \ \{:id :e-coherence, :relation :requires, \dots\ :wanted :naturality-and-coherence-for-monoidal-structure\}],\\
\ :holes [\{:kind :missing-warrant, :edge :e-monoidal-def, \dots\} \{\dots\}]\}
\end{minipage}
\end{center}

Two things are worth noticing, because they are the whole point of the IATC
layer. First, the edge from ``composites are semi-groupal'' to the monoidal
definition is typed \texttt{:defines} and carries an explicit
\texttt{missing-warrant} naming what the prose skipped:
\texttt{semi-groupal-} \texttt{composition-is-closed-} \texttt{under-composition}.
The word \emph{Clearly} in line 176 is exactly the deletion the warrant records.
Second, every such hole is mirrored in the top-level \texttt{:holes} vector,
which is what the gate checks; a graph asserting a hole it does not declare is
rejected.

\subsection*{A.3 What S7 made of it, and S4 beside it}

\begin{center}\small
\begin{minipage}{0.46\textwidth}\ttfamily\footnotesize\raggedright
\{:clean/proof "2311.05789\_\_p0"\\
\ :clean/typing-source :llama\\
\ :clean/seq [:construct-auxiliary-object\\
\ \ \ \ \ \ \ \ \ \ \ \ \ :reduce-to-known-result]\\
\ :clean/boxes\\
\ [\{:id :e-monoidal-def\\
\ \ \ :method :construct-auxiliary-object\\
\ \ \ :consumes [:semi-groupal-composition]\\
\ \ \ :produces :monoidal-category-def\\
\ \ \ :hole \{:kind :sorry :discharge :sorryProof\\
\ \ \ \ \ \ \ \ \ :satiety :payoff\\
\ \ \ \ \ \ \ \ \ :wanted "semi-groupal-composition-\\
\ \ \ \ \ \ \ \ \ \ \ \ \ \ \ \ is-closed-under-composition"\}\}\\
\ \ \{:id :e-coherence\\
\ \ \ :method :reduce-to-known-result \dots\}]\}
\end{minipage}\hfill
\begin{minipage}{0.46\textwidth}\ttfamily\footnotesize\raggedright
\% S4, leaf-0001 (L102--120)\\
\{:kind :connection/application-domain\\
\ :source \{:lines [102 120]\}\\
\ :slot-fill \{:application "condensed matter\\
\ \ \ physics (topological states of matter)"\}\}\\[6pt]
\% S4, leaf-0003 (L392--425)\\
\{:kind :connection/example-source\\
\ :source \{:lines [392 425]\}\\
\ :slot-fill \{:example "the category of sets,\\
\ \ \ abelian groups, and vector spaces over\\
\ \ \ a fixed field k"\}\}\\[6pt]
\% S4, leaf-0002 (L225--276)\\
\{:kind :connection\\
\ :slot-fill \{:other-structure "braided and\\
\ \ \ symmetric monoidal categories"\}\}
\end{minipage}
\end{center}

S7 types the two inference steps into the method vocabulary
(\texttt{construct-auxiliary-object} then \texttt{reduce-to-known-result}) and
carries the IATC hole through as a \texttt{:sorry} with a discharge obligation
--- which is the bridge to the APM formalisation machinery, the same
\texttt{:sorry} vocabulary on both sides. S4, working the same paper's prose
rather than its proofs, recovers what the introduction is \emph{doing}: naming
an application domain, offering an example source, connecting to a neighbouring
structure. Neither layer sees the other's spans.

\subsection*{A.4 The two open gaps, as fixtures}

The following are \textbf{illustrative fixtures, not measurements} --- they show
the shape the missing artifacts would take, so that closing \hz{H13} and
\hz{S11} is a matter of producing this rather than of deciding what to produce.

\paragraph{\hz{H13} --- what an S5 row looks like when rung-3 exists.}
Today every row reads as the left column: the strategy half is absent, so the
verdict cannot discriminate. With \texttt{rung3\_technique} promoted to a
stage, the same rows would carry the right column, and \texttt{weak-proof}
becomes reachable --- distinguishing \emph{we extracted badly} from \emph{the
proof is thin}, which is the distinction S5 exists to draw.

\begin{center}\small
\begin{tabular}{l|ccccc|l}
 & \texttt{noun} & \texttt{s:r3} & \texttt{s:pr} & \texttt{strat} & \texttt{comp} & \textbf{verdict}\\
\hline
\emph{today} (real) & 1.00 & \texttt{--} & 0.00 & 0.00 & 0.00 & \texttt{weak-extraction}\\
\emph{today} (real) & 1.00 & \texttt{--} & 0.00 & 0.00 & 0.00 & \texttt{weak-extraction}\\
\hline
\emph{fixture} & 1.00 & 0.75 & 0.40 & 0.62 & 0.81 & \texttt{comprehended}\\
\emph{fixture} & 1.00 & 0.20 & 0.10 & 0.15 & 0.34 & \texttt{weak-proof}\\
\emph{fixture} & 0.25 & 0.60 & 0.55 & 0.58 & 0.31 & \texttt{weak-extraction}\\
\end{tabular}
\end{center}

\paragraph{\hz{S11} --- what the structural canon looks like when it persists.}
S11 computes a canonical-shape census and whole-paper signatures and then
discards them. The fixture below is the artifact shape we propose:
\texttt{<run-dir>/structural-canon.json}, one record per canonical shape with
its population and exemplars, plus one signature per paper. The real S7 macro
census (5 archetypes, $67\%$ \emph{reduce-to-known}) is the evidence that a
census of this shape is computable; what is missing is only that S11 writes it.

\begin{center}\small
\begin{minipage}{0.86\textwidth}\ttfamily\footnotesize\raggedright
\{"run\_id": "mark7z", "corpus\_id": "math-ct-e2e-12",\\
\ "shapes": [\{"shape": "define-then-verify-coherence", "n": 34,\\
\ \ \ \ \ \ \ \ \ \ \ \ \ "exemplars": ["2311.05789\_\_p0", "0708.1921\_\_p2"]\},\\
\ \ \ \ \ \ \ \ \ \ \ \ \{"shape": "reduce-to-adjunction", "n": 21, "exemplars": [\dots]\}],\\
\ "signatures": [\{"paper": "2311.05789", "archetype": "definitional-survey",\\
\ \ \ \ \ \ \ \ \ \ \ \ \ \ \ "seq": [:construct-auxiliary-object :reduce-to-known-result]\}]\}
\end{minipage}
\end{center}

\subsection*{A.5 What this annex is evidence for, and what it is not}

It is evidence that the layers do what the warrant table says: S1 marks, S3
recovers a typed argument DAG with declared holes, S4 recovers rhetorical
structure from the same paper's prose, S7 types the argument into a method
vocabulary carrying the hole forward as a discharge obligation. All of it is
checkable against the source printed beside it.

It is \emph{not} evidence of corpus-scale quality: one paper, chosen for being
displayable. Its value is the opposite of the census's --- the census counts
many things approximately, the annex checks a few things exactly, and \hz{H21}
is precisely the defect that only the second can see.
